Neural tactic policies are routinely used to guide best-first proof search, but their benefit over a careful hand-written ordering is rarely measured in a setting where the baseline can be tuned and optimal proofs can be computed.
We build a tiny propositional sequent calculus (invertible G3cp-style rules, tactic = which compound formula to decompose), generate \rhval{const/n_train_seq} provable training sequents with 3--10 connectives, and train a small transformer and a feature-based MLP to imitate optimal tactics (minimum proof size by dynamic programming).
We then compare best-first search guided by the policies against breadth-first search (BFS) and a tuned hand heuristic (defer branching rules) within a budget of \rhval{const/budget} expansions on sequents with 11--20, 21--40 and 41--70 connectives, with \rhval{const/n_seeds} seeds, all systems sharing one search engine.
The result is largely negative: the transformer policy beats BFS in every test bin (e.g.\ solved rate \rhval{main/transformer-policy/test_41_70/solved/mean:2} vs.\ \rhval{main/bfs/test_41_70/solved/mean:2} at 41--70 connectives) but loses clearly to the hand heuristic (\rhval{main/hand-heuristic/test_41_70/solved/mean:2}), and the gap does not shrink with size.
Search with the policy's cumulative negative log-probability as cost is the weak point: following the same transformer greedily, or ranking tactics by the policy instead of using $-\log p$, raises the 41--70 solved rate to \rhval{abl_search/transformer-greedy/test_41_70/solved/mean:2}, on par with the heuristic.
A training-set sweep shows the opposite of the usual trend: the transformer trained on 1000 sequents extrapolates better (\rhval{analysis/analysis/test_41_70/nt1000_t41_m/mean:2} at 41--70) than the one trained on \rhval{const/n_train_seq} (\rhval{analysis/analysis/test_41_70/nt20000_t41_m/mean:2}), but still below the heuristic.
In a post-hoc (not pre-registered) variant, the feature MLP with rank cost reaches \rhval{abl_search/mlp-rank-cost/test_41_70/solved/mean:2} and exceeds the heuristic.
All conclusions are limited to this synthetic, small-scale setting.
